\section{Uniform comparison of initial environments}

Let $x=(T,u)$ be a finite initial tree with a tagged core vertex $u$.
Only core children born after time zero, and their subsequent descendants,
are assigned positive genealogical generations. Write
\[
 A_m(s;x)=
 \mathbb E_x\sum_{\substack{v\text{ born after }0\\
                          \mathrm{gen}_u(v)=m}}
 e^{-s\tau_v},\qquad m\ge1,\quad s>0.
\]
This is the scalar reward represented by $\mathcal K_s^m\mathbf1$
in Section~\ref{model:section}.
It is not an arbitrary function of the entire ambient tree.

A birth state means that the tag has just been promoted: $L_u=c_u=1$,
so its full degree is two. The original core neighbours of any initial
core vertex remain neighbours forever, and its full degree never
decreases.

\subsection{A common potential genealogy and the comparison identity}

Use a rate-one clock for each actual core vertex. A newly promoted leaf
can be given a fresh independent rate-one core clock, by the Markov
property. Introduce a potential genealogy rooted at $u$ in which
\emph{every} ring of a potential parent's clock creates a potential child,
whether or not that ring creates an actual core child. Label children by
their parent's ring index and run unused potential branches independently.
The actual future core genealogy is a subset of this common potential
genealogy. Its generation counts satisfy
\[
 \mathbb E\overline Z_j(t)=\frac{t^j}{j!}.
\]

Suppose two environments are coupled so that their descendant histories
agree until an increasing event $\Gamma(t)$ occurs. Suppose also that
$\Gamma$ depends only on clock streams of initial core vertices other
than $u$. These streams can be chosen independently of the potential
future genealogy. Each potential generation-$m$ birth has an inclusion
indicator in each model. Their difference has modulus at most one, and
is zero before $\Gamma$. The compensator of potential generation-$m$
births is $\overline Z_{m-1}(t)\,dt$. Hence
\begin{equation}\label{env:comparison}
 \boxed{
 |A_m(s;x)-A_m(s;x')|
 \le
 \int_0^\infty e^{-st}
 \frac{t^{m-1}}{(m-1)!}\,
 \mathbb P(\Gamma(t))\,dt .
 }
\end{equation}
There is no factor two: both actual genealogies are subsets of the
\emph{same} potential genealogy, and the difference of their two
inclusion indicators is bounded by one. Independence is invoked for the
potential event $\Gamma$, not for actual accepted arrows whose rates
depend on the evolving tree.

\subsection{Finite-radius comparison}

Decorate the initial core by the full initial labels $(L_i,c_i)$,
including the numbers of core neighbours outside any displayed ball.

\begin{theorem}[Uniform spatial locality]\label{env:locality}
Suppose two initial environments have isomorphic decorated radius-$R$
balls around their tags, where $R\ge1$. Let $c_u$ be the common initial
core degree of the tag. Then a coupling satisfies
\[
 \mathbb P(\Gamma_R(t))
 \le\min\left\{1,\frac{c_u t^R}{R!}\right\},
\]
and therefore
\begin{equation}\label{env:locality-integral}
 |A_m(s;x)-A_m(s;x')|
 \le
 \int_0^\infty e^{-st}\frac{t^{m-1}}{(m-1)!}
 \min\left\{1,\frac{c_u t^R}{R!}\right\}\,dt .
\end{equation}
For every $T>0$, a convenient explicit bound is
\begin{equation}\label{env:locality-tail}
 \frac{c_u T^{R+m}}{(R+m)R!(m-1)!}
 +s^{-m}e^{-sT}\sum_{j=0}^{m-1}\frac{(sT)^j}{j!}.
\end{equation}
For fixed $m,s$ and a birth state, this tends to zero exponentially in
$R$ on taking $T=aR$ with any fixed $0<a<1/e$.
\end{theorem}
\begin{proof}
For an initial core vertex $i$ with $c_i\ge1$, split its rate-one clock
by first selecting uniformly one of its $c_i$ initial core neighbours.
With probability $c_i/D_i(t)$ attach to that selected neighbour;
otherwise select uniformly among all other current neighbours.
This gives probability $1/D_i(t)$ to every current neighbour and hence
is an exact construction. The potential arrows along each directed
initial core edge $i\to j$ have rate $1/c_i$.

Allow arbitrary discrepancies at the radius-$R$ boundary from time zero.
To affect the tag, any external discrepancy must traverse the unique
initial-core path to the tag by inward arrows in chronological order.
New branches cannot bypass an initial-tree cut edge. Let $\Gamma_R(t)$
be the event that some such path has its required potential arrows
by time $t$.

For a path $u=v_0,v_1,\ldots,v_R$ directed outward from $u$, the expected
number of inward ordered arrow tuples is
\[
 \frac{t^R}{R!}\prod_{j=1}^{R}\frac1{c_{v_j}}.
\]
If $p(\gamma)=\prod_{j=0}^{R-1}c_{v_j}^{-1}$ is the probability that
simple random walk follows that outward path, then
\[
 \prod_{j=1}^{R}\frac1{c_{v_j}}
 =\frac{c_u}{c_{v_R}}p(\gamma)\le c_u p(\gamma).
\]
The probabilities of these length-$R$ outward paths sum to at most one.
A union bound proves the claimed estimate. All source clocks used by
$\Gamma_R$ belong to initial vertices other than $u$, so
\eqref{env:comparison} applies. Cutting the integral at $T$, bounding
$e^{-st}\le1$ on its first part, and evaluating its gamma tail gives
\eqref{env:locality-tail}.
If $c_u=0$, the core is a singleton and there is no exterior to compare.
\end{proof}

The boundary labels are necessary for exact agreement of the interior
initial states. The theorem is an approximation result; it does not say
that the radius-$R$ state itself is Markov.

\subsection{An exact limiting boundary for a high-degree parent}

The absorbing-parent process starts with the tag, its one leaf, and an
inert external core-neighbour stub. The stub contributes one to the tag's
degree. A ring selecting the stub has no effect on the retained process;
it does not kill that process. All remaining dynamics are unchanged.
Denote its scalar generation rewards by $A_m^\infty(s)$.

\begin{theorem}[Uniform parent-degree estimate]\label{env:parent}
If $x=(T,u)$ is a birth state whose parent has initial full degree $d$,
then
\begin{equation}\label{env:parent-bound}
 \boxed{
 |A_m(s;x)-A_m^\infty(s)|
 \le s^{-m}-(s+1/d)^{-m}
 \le\frac{m}{d\,s^{m+1}}.
 }
\end{equation}
\end{theorem}
\begin{proof}
Reserve a potential arrow from the parent to $u$ at constant rate $1/d$.
At such a ring, accept the arrow with probability $d/D_p(t)$.
All rejected and other rings are routed uniformly among the remaining
current neighbours. The accepted rate to $u$ is $1/D_p(t)$ and every
other current neighbour has the same rate, so this construction is exact.
The inequality $D_p(t)\ge d$ makes it valid at every time.

Until the first potential parent arrow, the tag's descendant histories in
the original and absorbing-parent processes agree. The potential arrow
stream is independent of the potential descendant genealogy, and its
probability of ringing by time $t$ is $1-e^{-t/d}$. Substitution in
\eqref{env:comparison} gives the difference of the two gamma integrals
in~\eqref{env:parent-bound}. The final inequality follows by integrating
the derivative of $x^{-m}$ from $s$ to $s+1/d$.
\end{proof}

The limiting process still contains interactions among the tag and its
new descendants. The theorem does not replace that process by independent
families or by a scalar birth-rate model.

\section{A finite approximation of the initial environment}

Fix a birth state, a depth $R\ge1$, and an integer $M\ge3$.
Traverse its initial core outward from the tag. Remove a component at
the first vertex on a path whose initial full degree is at least $M$,
or at depth $R$, whichever is encountered first. Replace each removed
component by an inert core-neighbour stub attached to its retained parent.
Retain all initial leaf counts and all core-neighbour counts of retained
vertices, including their stub neighbours. Write $\pi_{R,M}x$ for this
finite environment with boundary stubs.

Every retained vertex has depth less than $R$ and full degree less than
$M$. Thus, up to rooted isomorphism, the range of $\pi_{R,M}$ is a
finite set of decorated initial environments. Its future evolution is
still an infinite-state process.

\begin{theorem}[Uniform finite initial-environment reduction]
\label{env:finite-initial}
For every birth state $x$, every $m\ge1$, $s>0$, and $T>0$,
\begin{equation}\label{env:finite-bound}
\begin{split}
 |A_m(s;x)-A_m(s;\pi_{R,M}x)|
 \le s^{-m}\left\{
 \frac{T^R}{R!}+\frac{e^T-1}{M}
 +e^{-sT}\sum_{j=0}^{m-1}\frac{(sT)^j}{j!}
 \right\}.
\end{split}
\end{equation}
In particular, for fixed $m,s$, first choosing $T$, then $R,M$,
gives arbitrarily small uniform error over all finite initial environments.
\end{theorem}
\begin{proof}
The depth boundary contributes the potential-event bound $t^R/R!$
from Theorem~\ref{env:locality}, since a birth state has $c_u=1$.

For a first high-degree cut vertex $v$ at distance $r$, reserve its
parentward potential arrow at rate $1/D_v(0)$ and accept with probability
$D_v(0)/D_v(t)$, as in Theorem~\ref{env:parent}.
At intermediate vertices use the old-core potential arrows of rates
$1/c_i$. Cut vertices never occur as intermediate vertices on another
retained cut path, so these constructions are compatible.

If $\gamma=(u=v_0,\ldots,v_r=v)$ is the path to the cut vertex, the
product of inward potential rates is
\[
 \frac1{D_v(0)}\prod_{j=1}^{r-1}\frac1{c_{v_j}}
 =\frac{c_u}{D_v(0)}p(\gamma)
 \le\frac{p(\gamma)}{M}.
\]
First-cut paths are prefix-free, so their simple-random-walk cylinder
probabilities sum to at most one, even when their lengths differ.
The probability of some high-degree influence reaching $u$ by $t$ is
therefore at most
\[
 \frac1M\sup_{r\ge1}\frac{t^r}{r!}
 \le\frac{e^t-1}{M}.
\]
Combining the two kinds of boundary gives a potential event with
probability at most
$\min\{1,t^R/R!+(e^t-1)/M\}$.
Use \eqref{env:comparison}; on $[0,T]$ take its increasing probability
bound at $T$, and use $\int_0^\infty e^{-st}t^{m-1}/(m-1)!\,dt=s^{-m}$.
On $(T,\infty)$ use the probability bound one and the exact gamma tail.
This proves~\eqref{env:finite-bound}.
\end{proof}

The finite range concerns initial environments. Each retained profile still evolves on an infinite state space; the population truncation below supplies a finite calculation at every fixed time.

